% Extended lesson notes for the final group of problems.

\expandafter\def\csname sourceexcerpt@50\endcsname{%
altogether and so the data structure algorithm we're going to use is called the prefix
and suffix product arrays technique so prefix and suffix product arrays technique and
some may Wonder initially and I thought of this myself what about two pointer I've
heard of this thing called two pointer but isn't that the same as a prefix and suffix
product arrays technique well no it's not you see the differences between two pointers
is that we're focusing on really just the whole array all of itself and we're going to
be focusing on the left and right side of the array not particularly two points in the
array we are going to be scanning and creating two different arrays to be able to keep
track of what we have already done so to answer what are we going to do with the data
we will we are going to first create two arrays which will then be vectors in C++ to
store all the elements to store all the elements visited create one array in fact I'm
just going to put this inside of the technique so that way we know exactly what to do
right over here so we would first create two arrays to store all the elements visited
create one array for the left side visited and then create one array for the right
side visited and then finally what are we going to do with the data we will loop
through our given array create both the left and the right side arrays by creating a
formula to take the index of each element and store them in their respective arrays
then multiply both the right and left side arrays we created for the output.
}

\expandafter\def\csname sourceexcerpt@51\endcsname{%
us to do just run the problem is that we are given a head of a linked list and the it
looks like here okay and wants us to remove the nth node from the end of the length
list okay so from the end of the linked list so we have the second note accordingly and
we remove that and therefore we have one 2 three and five same thing for well if we
have one in particular and we remove the first index since we are removing from the
end of the list that's tricky to start we're then going to have one accordingly so
that's why in particular we have this input and the output so from the end of the link
list the output is a new linked list with the output or hang on with the nth node
removed. The technique that helps solve this problem is two pointers. We have
two elements of data we will put at the head of the linked list as well as a as well
as a second pointer after the first pointer and this will help us scan throughout the
rest of our linked list given the fact that we were just given the head so we can remove
our nth node okay and the following steps in order to break down the following problem
are as followed so for step one we're going to initialize our two pointers at the head
of the list then we will move the first pointer n nodes ahead from the end and we'll
use a for Loop to iterate through and.
}

\expandafter\def\csname sourceexcerpt@52\endcsname{%
in the list nodes only nodes themselves can be changed okay so we need to focus on the
nodes position so we need to change the nodes position based on their desired order
and we need to return a modified linked list all right so the technique for us to be
able to solve the following problem that we would need in order to rearrange the order
is that we are going to be using a fast and slow pointer method it's the tortois and
her algorithm where we have a fast and slow pointer and we focus on dividing the
linked lists based on their halves reverse one of the halves then merge them together
so this steps for us to be able to solve the following problem or as follows step one
Hang on we're going to have the edge case where we have a null pointer step two we're
going to initialize two pointers at the head of the list to find the middle in step
three we're going to move the first pointer to two steps at a time and the second
pointer one step at a time then the next part we're going to reverse the second half
of the list the fifth step we're going to merge the two halves taking one node from
each half alternatively then we're going to initialize our two halves and expand upon
the relationship between the first and second nodes and finally we're going to set the
second node as the next node so to give the following what we're going to do our
following steps I'm just going to move all this up here so have step one step two step
three step four step five is.
}

\expandafter\def\csname sourceexcerpt@53\endcsname{%
we are given putting up our questions here our answers down here skop input so right
here we are given a linked list with a head node hence right here and then the data
structure that we would need or Technique we have to use we would use an iterator
approach with three pointers to reverse the linked list want to do with the data we will
modify pointers of each node to reverse the direction of the links and then for output
we will return the new head of the reversed linked list and so now this is just at the
point of myself showing you how to do it so we will then just go over the steps over
what we need to do which will break it down into five steps so putting this in our
section one we need to initialize three pointers which will be the previous where we
will initialize as null because everything is connected current as head and next as
null because we haven't seen anything yet all right and we're going to Loop through
the linked list until current becomes null since we're moving up we're treating this
like an array third we each iteration we will assign next as the next node of current
so technically we have four pointers to start change the next pointer of current to
previous and reversing the link move pre to current and current to next finally after
the loop we will assign the head as the new head which is now the last node of the
original linked list so just rundown linked list is pretty easy to understand it's like an
array but instead it's connected that's just the best way of.
}

\expandafter\def\csname sourceexcerpt@54\endcsname{%
the main technique is bit manipulation and I'll also have some signs with some
operators as well that we use for bit manipulation in order to solve a following
problem but what we will be doing is that if we have a 32 bit unsigned integer that we
have we just need to reverse that and finally have our numerical output at the same
time so as indicated over here and this requires the use of binary operators to be
able to manipulate our far our following binary input so what we need to do so we need
to return the Reversed bits there's some notes on here about Java but we're focusing
on C++ what to do with the data reverse the bits and then input we are given a 32 bits
unsigned integer and there is a key difference between signed and unsigned integers
in particular I also have another visual to also further explain that will be up here
and without further Ado let's be able to write out the solution for this because it's
going to be pretty straightforward over exactly we what we need to do this will be
taken in three steps so step one what we need to do is that we need to in initialize
the result and second step we need to do is initialize a variable to keep track of the
number of bits to be processed in step three we need to reverse the bits of the given
integer using a while loop and for the first part of three inside the while loop we're
going to shift the result to the left by one and then we are going to shift the given
integer to the.
}

\expandafter\def\csname sourceexcerpt@55\endcsname{%
in the parameters to be able to see exactly what's going to be input it is represented
an image rotate the image by 90° okay clockwise you have to rotate the image in place
which means you can modify the input 2D Matrix directly do not allocate another 2D
Matrix and do the rotation okay so that means we need to modify and that means we're
not going to be necessarily storing per se okay we won't be storing anything since
right here if we are rotating an image in place we just need to modify the 2D Matrix
directly and that means we're not going to make another Matrix to store every single
element that we are visiting this part can be tricky at first but if you understand
more about data structures and algorithms data structures can be used in particular to
keep track of every single element that we have visited in order for our output which
would be our answer however we are not going to be using a data structure in this case
instead we are just going to be doing this technique called an inplace Matrix rotation
and what is that really well an in place Matrix rotation is really just two for Loops
where we swap both of the first element we scan and the second element we also scan as
well so for instance if we were to have a matrix of 1 2 3 4 5 6 789 what we need to do
really in terms of rotating it you notice the n the numbers one two and three they're
all going to be rotated 90 degrees so 1 2 and three accordingly and then also the same
thing 456 which.
}

\expandafter\def\csname sourceexcerpt@56\endcsname{%
checking to see if they are the same or not and then we will be returning true if the
same so the technique we will be doing is depth-first search namely to be able to make a
comparison between the two where depth-first search pretty works and is really extensive
given its space and time complexity that we would use and keep everything within the
constant space and so what we will be doing is that we're going to be checking on the
different values between all of our n u all between our pointers between if one or one
of P or Q is null or the values are not equal to each other we're going to be having
one condition where it's going to be false and then we're going to see in particular
if prq are going to be matching up so use this technique to check their pointers first
and then we also do a recursive call to check to see the left and right sub trees of p
and Q just from the roots for def search is known for recursion so to be able to break
down the steps or what we are going to do step one we are going to we're going to be
checking if p and Q are null return true and then for the second part if one of P or Q
is null or their values are not equal return false and then we're going to be
recursively calling the function for the left sub trees of p and Q and then we will be
doing the same thing but for the left sub trees so the four boundaries so we're doing this for
left and the right.
}

\expandafter\def\csname sourceexcerpt@57\endcsname{%
string and this string can be deserialized to the original tree structure. Serialization also
the input SL output format is the same as how LeetCode displays a binary tree you do not
necessarily need to follow this format so please be creative and come with a different
approach yourself okay so for what we will be doing for our following steps is that
for one we are given the root of a binary tree as our input so I'm just going to label
out the rest of our following questions over what we need to do so we know exactly
what input is the root of binary tree what data structure algorithm technique to use
what to do with the data and finally what our output is going to be so we're given the
root of binary tree it looks like from our input to Output or output for serialization
we need a string that encodes the structure of the tree for deserialization the root
node of the reconstructed tree so then the technique that we're going to be doing for
this since we're focusing connectivity of a given trees that we're going to do depth
first search using recursion for both serialization and d serialization and so the
steps for us to be able to go about for each one are going to be The Following over
here for two functions I'm going to put these notes in the comments as well is that
Force heroization we will overall we will Traverse the tree in preorder traversal and
build a string representation using essential value for null noes for D serialization
we use preorder string representation to reconstruct the tree creating new nodes for
each value and null for.
}

\expandafter\def\csname sourceexcerpt@58\endcsname{%
that row and column to zero then we are going to iterate again if the first cell of a
row or column is zero we set the entire row or column to zero accordingly so that
means we're going to be looping throughout the rest of our given throughout the rest
of our given Matrix that we will have right here let's put the notes down here zero
right here okay and given empo that we have we have a Matrix and we are supposed to
return a matrix with corresponding zeros so what we will be doing is that for a
following way we are going to first 2D Matrix so the steps we are going to be doing to
be able to get through all of this is that for the first step if a zero is found at
cell let's just say for example cell I and cell I of J as our initial one we're going
to set the Matrix accordly this one and Matrix Z J to Zer accordingly step two we're
going to iterate over the Matrix again to indicate we have a modification in place if
Matrix the four boundaries I of zero or Matrix whenever we have a row or column is zero we're
going to set Matrix I n j to zero then step three we're going to be outputting the
modified Matrix and so I'm going to be expanding upon the pseudo code to be able to
help break down over what we need to do so what we will be doing in this case is that
for each cell in The Matrix for step one if the cell value is zero we're going to set
the first cell of the corresponding.
}

\expandafter\def\csname sourceexcerpt@59\endcsname{%
elements are returned in spiral order. The solution uses an iterative approach with
four boundaries: top, bottom, left, and right. These boundaries control each direction
of the scan
so this may be a little bit confusing at first but I'm going to be explaining more
within the steps that we have and so we will have in order for us to be able to
replicate this entire behavior of the iterative approach what we need to do is that
first we need to Define our four pointers top bottom left and right we also initialize
list which will be Vector since naturally arrays and list can be treated as vectors
and such and so the next step that we will do is that we are going to run a while loop
from Loop until the top of the Matrix is less than or equal to the bottom and the left
of the Matrix is going to be less than or equal to the right and what we're going to
do is that we're going to print the top row from left to right and then increment the
top pointer and then we will print the right column from top to bottom decrementing
the pointer right pointer and then we will be reaching to this point if the top is
less than or equal to the bottom we're going to print the bottom row from right to
left and then decrement the bottom pointer and then if the left pointer is less than.
}

\expandafter\def\csname sourceexcerpt@60\endcsname{%
so if we have the root of one binary tree and then the root of a sub of a sub tree
accordingly of sub root and then a root we're checking to see if these values are
matching so that's all we are doing the values from the sub root are matching any of
the sub trees from the root of the binary tree it's pretty self-explanatory and then
we return true or false just seeing if these nodes accordingly within the root of our
B inary tree have the same values as the following right over here so just to be able
to break down even further the technique that we will be able to use to be able to
solve this problem is that we're going to be doing a recursive approach to check if
one tree is a sub tree of another and what we will be doing with our data aside from
checking we will Traverse the trees recursively comparing the nodes and their sub
trees so then what we will be doing next I'm just going to be going over the steps
over how we will be solving this problem and let's just double check for case of that
yep it's pretty self-explanatory between examples one and two check and see if these
sub trees have accordingly the same type of numbers already between sub tree as well
as the binary tree now this example may have the sub root values but then we have this
zero at the end that will also expand upon two so if we were to break down how a
binary tree works for the most part it typically will start with the root then you go
left then you.
}

\expandafter\def\csname sourceexcerpt@61\endcsname{%
then we initiated for instance A + B and then return sum this is not what they're
going to look for watch this never mind this works strangely enough we can do that
just kidding I wish just a little humor to help you guys out during the whole time but
anyway I'm just going to give you guys the steps over exactly what to do on this whole
thing so steps to be able to go through on this so step one what we need to do is that
we need to initialize our two integers then we also need to initialize an unassigned
integer we're going to call this one carry then the third step of what we're going to
do is that we are going to perform bit one operations to calculate the sum step four
is that while one of the numbers can be a or b that we initialized doesn't equal zero
we calculate the sum by using the carry zor binary operator then what we will do is
that we're going to calculate the carry by using and operator with left shifting next
part we're going to set the first integer to be the sum then we will set the second
integer to be the carry and finally we will return the first integer so to be able to
go over the steps what to do so step one step two step three step four step five on
what to do step one we're going to ini size everything along with step two at the same
time so int num of one is going to be equal to a just for clarity int num of two is
going to be equ Al to B.
}

\expandafter\def\csname sourceexcerpt@62\endcsname{%
what we're working with so we have a given array of nums and the K most frequent
elements it looks like it's just asking us to return the two most frequent elements so
we have 1 two three ones we have three ones accordingly then we have it looks like two
twos and just 1 three so there are three ones and two twos the two most frequent
elements are yes one and two and it looks like that we need to return this in an array
as they are expected and it looks like this Edge case if it's just one and the one
most frequent element it will just be one also some more constraints as well okaye the
range one the number of unique elements inray so it starts at one and follow up your
algorithms time Collex must run better than o n log n where n is the arrays size okay
so that's also another thing to take into consideration and then the data structure
and technique that we use is really just going to be a hash
map with a priority queue and to explain the data structures between the two over exactly
how we're using them is that a hash map is much like in the case of a series of
lockers that one would have if one has a key that needs to connect to a specific
Locker the value you have a key value pair to be able to open your locker and that is
how the hash map is typically designed and then the priority Cube for the most part is
that cues in general are much like in the case of grocery store lines where.
}

\expandafter\def\csname sourceexcerpt@63\endcsname{%
minus1 nus one the robot can only move either down or right at any point in time given
the two integers M and N return the number of possible unique paths that the robot can
take to reach the bottom right corner okay so we need to return all the possible pass
input we are given two integers which represent a grid as well okay so the technique
that we're going to be doing is going to be dynamic programming where we're just going
to expand upon the variables and their relationships and so we need to ask ourselves
what to do with the data so the first step we need to do is that we are going to initi
we are going to initialize a 2d grid which will be a vector inside of another Vector
of integers called DP and the Second Step we're going to do is set the value of the
first row and First Column cells to one since there is one cell and this is how we set
everything up in the third we're going to Loop over the grid starting from the second
row and second column then for each cell going to set its value equal to the sum of
the cell above it and the cell to the left of it and we are just going to finally
return the number of unique paths which would be represented as our red that we have
created M minus one n minus one as the exact values okay so first things first going
do step one which will be Vector here Vector int DP m Vector int then n we're going to
do step two initialize our Rows 4 into I is going.
}

\expandafter\def\csname sourceexcerpt@64\endcsname{%
letters exactly once so in our first example we have the input of s is equal to n a g
and T is equal to n a gr whenever we rearrange the words all together will be the same
and so we are we already found that these are the same word and so it's considered
true are second example has the input of rat and car though we have the same number of
characters these are totally different words in particular because they're not the
same characters and so we also have our constraints S\&T consists of lowercase English
letters that's good and so to be able to give a breakdown from our problem we already
know that we have two strings SNT and the following data structure algorithm technique
we're going to use we're going to use one map to keep track of the frequency of
characters for both s andt and what do we do with the data it's going to be the
following steps overall we will keep track of each character in our maps to determine
if they both have the same number of characters and the same type so the steps that we
need to do is that first we need to initialize our map to keep track of the frequency
of similar characters between variables s and t then second step we are going to Loop
through the S string and increment the frequency of characters for S then next part we
will Loop through each character in string T and doing so we will decrement each
character and the frequency map and overall finally if there is a nonzero number we
will return false if the frequency map didn't turn to zero.
}

\expandafter\def\csname sourceexcerpt@65\endcsname{%
I'll break it down and it looks like we will return a true or false condition and
Bings are set to True default false otherwise well if we set a condition turn all
right so A Man A Plan A Canal Panama true such as like okay so it's like damn it I'm
mad spell backwards say racecar there you go yay racecar spell
backwards is damn it I mad oh damn that's a p drone for the most part okay so that's
the same thing we got to check to see if it is a palad Drome or if it's not a palad
Drome because all the letters okay race a car and this is not a power drone because it
needs to be the same type of word forward and backwards all right that's the four boundaries same
space that is also a pal androme just like a number example three as well all right
and so the technique we're going to be using this time is going to be two pointers in
this case where it's not for two pointers you have one point at the end you have your
left and your right and it's like binary search but it's not binary search because two
pointers is not Divine and Conquer and you don't have a middle type of variable
whatsoever you have just have two pointers where it's just going to be your left and
your right of your string it looks like what we're going to be doing for steps is that
we're going to initialize our two pointers left and right then what we're going to do
is Loop until we're going to be expanding upon the relation ship.
}

\expandafter\def\csname sourceexcerpt@66\endcsname{%
we have a following input right here that would be if it's a parenthesis we're going
to return true and then if in example two all the parentheses are closed and they're
all corresponding so it's also true but in this case in the third one it's not they're
not the same brackets so it's going to be false accordingly so it looks like all
together we are given a string and we need to return true if valid parenthesis false
otherwise s consists of parenthesis only dot and we also need to false otherwise and
we also need to check for edge cases if there is no string as well no characters and
the string as well yeah Okie doie right here so now the technique I will be using is
we'll be calling upon a stack will and if you are curious and wondering what a stack
is it's a type of data structure that much like in a stack of plates that have various
methods of push pop and then such as push and pop on a stack namely one is able to
have organize their data accordingly depending on exactly how they want to keep track
of what has been scanned so we would use a stack to keep track of all the characters
visited and what are we going to do with the data well first things we're going to do
is that we're going to initialize our stack Second Step we're going to do is Loop
through the string for every character and then for all of the beginning parentheses
we will store the characters in the stack and then for the closing parenthesis we're
going to check to see if the top of.
}

\expandafter\def\csname sourceexcerpt@67\endcsname{%
so just by nature of binary search trees nodes that are going to be less than the noes
key or the noes value that we have with the current one it's going to be on the left
and then greater is on the right that's just based on principle when your search trees
by principal left is less right is more that's just what they mean and we just need to
replicate this entire behavior and nature of the binary search tree just to determine
if it's true and it looks like we are just given the root of a binary tree and then
our output is that return true or false if we have a binary a valid binary search tree
or not so the technique that we will be doing for the following case is going to be
recursion we will create a helper function to determine the current node the minimum
allowed value and the maximum allowed value and to go over the steps of how we will
replicate this one we're going to we're going to call upon a helper function we will
create and then call upon throughout the rest of our solution is valid binary search
tree with the root node and the minimum and maximum allowed values those are going to
be some we're going to create and then step two we're going to create our helper
function helper cursive function that takes the current node the minimum and Max
values then we're going to have a base case if the current node is null we're going to
going to return true and the reason being and even in even in the case of an empty
subtree it is also considered true then.
}

\expandafter\def\csname sourceexcerpt@68\endcsname{%
we have LeetCode being true just for this example and the reason being because return
true because LeetCode can be segmented as LeetCode but in this case apple and apple
and word Apple pen Apple ah okay so we have apple pen and apple pen it's like the one
example and then we return true and then return true because Apple pen Apple can
be segmented as Apple pen Apple note that you are allowed to reuse a dictionary word
okay and then we have cats and dog with a word dick of cats dog sand and Cat it's
false namely because hey we have cats right here and then we are just with left with
OG the it looks like all the words accordingly still need to be intact we can reuse
them but all the words need to be stay intact so just give wrak down our input we are
given a word s and a dictionary of strings called word deck and it looks like our
output is either true or if we can find the word in our word dictionary AR or false
otherwise all right and so the technique to be able to solve this following problem is
that we're going to use dynamic programming where dynamic programming best way of
describing it is what we will be doing is that we will be going forward and backward
Within A Word Dictionary to be able to find every single character that we find to be
able to make up our following string that we were given for Apple pen Apple we're
going to have to go forward then we also have to go backward to be able to find
exactly each word that will.
}

\expandafter\def\csname sourceexcerpt@69\endcsname{%
the words with the given Board of characters to find which words are in the board
where the characters are adjacent to each other and so the steps for us to be able to
accomplish this for the following test are going to be the following step one we will
initialize a tri data structure and insert all the words into it orang into it with
each character step two we will create a list to store the found words from our
character board step three we will Loop over each cell in the board using def for
search for both the row and the column step four you will return the list of found
words so just going over all the steps that we need to do I'm also going to further
expand this at the same time so step one what we will do ultimately is that to
initialize a tried data structure we will need to do the following create a separate
structure to initialize the trode children and word so if I'm going step one step two
step three step four step one is going to look like the following we're going to move
this out here so that way we can call upon a public struck that we have where it's
going to represent trodde and then we're going to have a vector or trodes for children
and then we're going to have a string called word then we're going to have a TR node
structur with children ranging from the English letters in the alphabet with a null
pointer and then word this is just how we initialize our following don't forget the
comments as well and this is just how we set up.
}

\expandafter\def\csname sourceexcerpt@70\endcsname{%
example two we have a given Matrix over here and then we have the word c is it all
connected well according to this one yes it's true now if we have this one a b c b as
this word within the board is it really connected all together well in this case it's
not so what we need to do is we need to replicate this Behavior to be able to have our
input or output exist accordingly so we would need to return true if the word is in
The Matrix accordingly and the technique we will be using for this is going to be
depth first search with recursion to check for each Matrix that has been checked
already and so I'm going to break down for search in the following stuff that we have
for what we are going to do with the following data so to be able to break down I'm
going to first expand in three different steps so one what we need to do is that for
our given solution we need to Loop over each cell in the board from top to bottom and
right to left accordingly using a nested for Loop and then for step two what we're
going to do we're going to perform depth-first search by creating a function which takes
in the parameters of the board index of I index of J with our word and our current
index therefore be returning true if the word is in The Matrix three otherwise return
false in this case and I'm going to be expanding a little bit more on part two namely
to be able to break down our following def first search function so.
}

\expandafter\def\csname sourceexcerpt@71\endcsname{%
need to write this algorithm in O login time complexity and overall just from looking
from the example we have nums equal to 45 67012 with the target of zero and our output
is four what we need to do is that we need to reflect the relationship of the
following examples example 1 2 and 3 to have our input reflect our output so if we
have a target of zero already started it's located right here so a Rays start at zero
by principle so 0 1 2 3 and 4 Target is zero right over here and it's located right
over here so it's at index 4 which we output four and if the target is not located in
the array we're going to have to return negative 1 somewhere in the case of right here
and over here and here so we're so we will use a technique
called binary search where in it's more of a divide and conquer method where we have
three variables of the middle low and high end and the way how we are going to figure
out between the two is that with binary search the principles the principle is say for
example you have a row of cards that's the best way of describing it and you are
supposed to find your given Target you have a value you need to find of one now you
don't know where one is already located because it's already been shuffled in the
amount of cards but what you can do is that if you have your following numbers already
taken into consideration for the array you can find where the middle is going to be
located and.
}

\expandafter\def\csname sourceexcerpt@72\endcsname{%
the length of longest strictly increasing sub sequence okay so we will return the
length of the longest strict the increasing subsequence we're given an array nums
right here so we need to figure out the streak of how many numbers it's pretty simple
all right to figure out the longest streak of numbers which are increasing it's pretty
self-explanatory we have here the longest increasing subsequence is over here four so
I'll just need to clarify for four ah okay so longest form of increasing subsequence
we have 01 and then 2 three 0 1 2 3 for the entire thing that would explain four we
have four elements over here so that makes sense just trying to clarify on that part
okay and but the tag we're going to be focusing on this problem, we will use
dynamic programming because we will be moving forward and backwards in our array to
understand the pattern and get our desired output most like going to have two Loops to
start with so going over the first St U we need to initialize an array called DP of
nums size we will for clarity decare and N is equal to num.

size three we will Loop over the array starting from the second number since we need a
sequence we will iterate also every number four the second number to have our sequence
e we will declare our second number to be greater than our first number yeah this will
be statement if our second number is greater than our first number we will set our
current index to get the max of number and DP and last member stored in our Ray of the
previous number + one.
}
